% !TeX program = xelatex
\documentclass[11pt,a4paper]{article}

\usepackage[UTF8,scheme=plain,fontset=fandol]{ctex}
\usepackage[margin=22mm]{geometry}
\usepackage{amsmath,amssymb}
\usepackage[hidelinks]{hyperref}
\usepackage{needspace}

\newcommand{\E}{\mathbb{E}}
\DeclareMathOperator{\Var}{Var}
\newcommand{\KL}{D_{\mathrm{KL}}}
\newcommand{\pref}{p_{\mathrm{ref}}}
\setlength{\parindent}{0pt}
\setlength{\parskip}{5pt}
\setlength{\emergencystretch}{2em}
\allowdisplaybreaks[1]

\hypersetup{
  pdftitle={Assignment 5: Language Model Fine-tuning},
  pdfauthor={Yilin Zhang},
  pdfsubject={Solutions to Assignment 5 of Generative Models}
}

\begin{document}
\setlength{\abovedisplayskip}{7pt}
\setlength{\belowdisplayskip}{7pt}
\setlength{\abovedisplayshortskip}{4pt}
\setlength{\belowdisplayshortskip}{5pt}

\begin{center}
  {\Large\bfseries Assignment 5: Language Model Fine-tuning}\\[8pt]
  Name: 张逸林
  \hspace{8mm}
  Student ID: 2500013576
\end{center}

\section*{1.}

\textbf{Pre-training.}
Pre-training uses large-scale text data, including news, papers, books, and code.
Its objective is to predict the next token in an autoregressive model.
The supervision comes from the next token actually present in the text, so this
is self-supervised learning, typically trained by minimizing cross-entropy loss.
Pre-training enables the model to learn general language patterns, knowledge,
and reusable representations, and provides initial parameters for fine-tuning.

\textbf{Instruction fine-tuning.}
Instruction fine-tuning uses pairs of instructions and reference responses, such as QA
interactions. Starting from the pre-trained parameters, the model is trained
with supervised learning to predict the tokens in the reference responses and
minimize cross-entropy loss. This teaches the model to follow instructions,
fit the chat format, and apply its existing knowledge to specific tasks.
It can be viewed as an extension of pre-training in the sense that it continues
to use a next-token prediction loss, but now trains on task examples.

\textbf{Reinforcement fine-tuning.}
Reinforcement fine-tuning uses prompts and model-generated responses, with
rewards assigned to evaluate the responses. These rewards may come from human
feedback, a learned reward model, or rules. The model updates its parameters
to maximize the expected reward of its responses, with the aim of better aligning
its outputs with user preferences and practical task requirements.

\textbf{Reference responses versus rewards.}
Learning from reference responses gives the model concrete examples of how to
answer, with a target token at each position. Learning from rewards instead
lets the model generate a response first and then uses a score to evaluate
the whole response. The score does not necessarily specify the correct next
token or a particular reference response to imitate. Different valid responses
may receive high rewards.

\section*{2.}

Since \(p(s;\theta)>0\) and \(p(s;\theta)\) is differentiable in \(\theta\),
the chain rule gives
\[
\nabla_\theta\log p(s;\theta)
=\frac{1}{p(s;\theta)}\nabla_\theta p(s;\theta).
\]
Multiplying both sides by \(p(s;\theta)\), we obtain
\[
\nabla_\theta p(s;\theta)
=p(s;\theta)\nabla_\theta\log p(s;\theta).
\]

By assumption, differentiation and summation may be interchanged, and
\(R(s)\) does not depend on \(\theta\). Therefore,
\[
\begin{aligned}
\nabla_\theta J(\theta)
&=\nabla_\theta\E_{s\sim p(\cdot;\theta)}[R(s)]\\
&=\nabla_\theta\sum_s p(s;\theta)R(s)\\
&=\sum_s R(s)\nabla_\theta p(s;\theta)\\
&=\sum_s R(s)p(s;\theta)\nabla_\theta\log p(s;\theta)\\
&=\E_{s\sim p(\cdot;\theta)}
  \left[R(s)\nabla_\theta\log p(s;\theta)\right].
\end{aligned}
\]

\newpage
\section*{3.}

\textbf{(a)}
Since \(R(s)=1\) for every response and \(\sum_s p(s;\theta)=1\),
\[
J(\theta)=\sum_s p(s;\theta)R(s)=\sum_s p(s;\theta)=1.
\]
Thus, \(J\) is constant, and every \(\theta\in\mathbb{R}\) is a global maximizer.

\textbf{(b)}
Since \(\eta>0\) and \(\frac{d}{d\theta}\log p(s;\theta)\neq 0\),
\[
\Delta\theta=\eta\frac{d}{d\theta}\log p(s;\theta)\neq 0.
\]
By part (a), the current \(\theta\) is already optimal, even though this
single-sample update is nonzero.

\textbf{(c)}
With \(R(s)=1\) and \(b=1\), every sampled response gives
\[
\Delta\theta_b
=\eta(1-1)\frac{d}{d\theta}\log p(s;\theta)=0.
\]
Hence, \(\Delta\theta_b\) is identically zero, and
\(\Var(\Delta\theta_b)=0\).

\Needspace{12\baselineskip}
\section*{4.}

Write
\[
Z=\sum_{s\in\mathcal{S}}\pref(s)e^{R(s)/\beta},
\qquad
p^\star(s)=\frac{1}{Z}\pref(s)e^{R(s)/\beta}.
\]
Since \(\mathcal{S}\) is finite, \(R(s)\in\mathbb{R}\),
\(\pref(s)>0\), and \(\beta>0\), we have \(0<Z<\infty\).
Also, \(p^\star(s)>0\) and \(\sum_{s\in\mathcal{S}}p^\star(s)=Z/Z=1\),
so \(p^\star\) is a valid probability distribution. All sums below are over
\(s\in\mathcal{S}\), with zero-probability terms interpreted as zero.

\textbf{(a)}
Taking the logarithm of \(p^\star(s)\) gives
\[
\log p^\star(s)=\log\pref(s)+\frac{R(s)}{\beta}-\log Z.
\]
Expanding the KL divergence, we obtain
\[
\begin{aligned}
\KL(p\Vert p^\star)
&=\sum_s p(s)\log\frac{p(s)}{p^\star(s)}\\
&=\sum_s p(s)
  \left[\log\frac{p(s)}{\pref(s)}-\frac{R(s)}{\beta}+\log Z\right]\\
&=\KL(p\Vert\pref)
  -\frac{1}{\beta}\E_{s\sim p}[R(s)]+\log Z,
\end{aligned}
\]
where the last step uses \(\sum_s p(s)=1\). Rearranging yields
\[
F(p)=\E_{s\sim p}[R(s)]-\beta\KL(p\Vert\pref)
=\beta\log Z-\beta\KL(p\Vert p^\star).
\]

\textbf{(b)}
By the non-negativity of KL divergence and \(\beta>0\),
\[
F(p)=\beta\log Z-\beta\KL(p\Vert p^\star)\leq\beta\log Z.
\]
At \(p=p^\star\), we have \(\KL(p^\star\Vert p^\star)=0\), so
\(F(p^\star)=\beta\log Z\). Equality holds if and only if
\(\KL(p\Vert p^\star)=0\), which is equivalent to \(p=p^\star\).
Therefore, \(p^\star\) is the unique global maximizer of \(F\), with
maximum value \(\beta\log Z\).

\end{document}
